\label{sec:resultados}

Este capítulo apresenta as medições da bateria descrita na
Seção~\ref{sec:analise}. Cada figura é uma grade de nove painéis. As linhas
trazem os dois modelos abertos, que instanciam o sistema inteiro (filtro e
componente principal), e o filtro fine-tunado da família gpt-oss, que ocupa só
o papel de filtro. As colunas trazem os três benchmarks. No eixo horizontal de
cada painel ficam as cinco configurações: sem filtro, \textit{prompt} simples,
delimitação, marcação e a combinação das duas. O NotInject tem apenas itens
benignos, então nele a taxa de sucesso do ataque não é definida e o painel
correspondente fica vazio.

Todo valor é a média das dez repetições da bateria, e toda haste vertical é o
intervalo de confiança de 95\% descrito na Seção~\ref{sec:protocolo}. Nas
figuras de valor absoluto, a haste mostra a incerteza de cada taxa; a barra
sem haste é uma taxa em que menos de cinco itens variaram, quase constante, e
sem intervalo pela regra da Seção~\ref{sec:protocolo}. Nas
figuras de diferença, cada diferença é um ponto: se a haste não cruza o zero,
a diferença não se explica pelo acaso; se cruza, falta evidência de efeito
nesta amostra, o que não equivale a ausência de efeito. As medidas de
validade e de custo operacional (saída inválida, resposta vazia, alcance da
marcação, \textit{tokens} e tempo) descrevem o comportamento observado e não
trazem intervalo.

\section{Taxa de sucesso do ataque}\label{sec:res_asr}

A Figura~\ref{fig:asr} apresenta a taxa de sucesso do ataque em valor absoluto.

% [NÚCLEO]
\begin{figure}[H]
\centering
\includegraphics[width=\textwidth]{asr}
\caption{Taxa de sucesso do ataque por configuração. O painel do NotInject fica vazio porque o conjunto não tem itens atacados.}
\label{fig:asr}
\end{figure}

A proteção total do \textit{guardrail} é a distância entre cada configuração e
a barra sem filtro. O filtro fine-tunado protege menos que o modelo aberto da
mesma família em todas as configurações; a Seção~\ref{sec:conc_guards} discute
o motivo.

A Figura~\ref{fig:delta-asr-simple} isola o efeito de separar o canal de
instrução do canal de dado não confiável. Cada condição é comparada ao
\textit{prompt} simples, com o mesmo filtro, os mesmos itens e o mesmo
componente principal.

% [NÚCLEO]
\begin{figure}[H]
\centering
\includegraphics[width=\textwidth]{delta_asr_vs_simple}
\caption{Diferença na taxa de sucesso do ataque de cada condição em relação ao \textit{prompt} simples. As hastes são intervalos de confiança de 95\% por \textit{bootstrap} BCa pareado sobre os itens.}
\label{fig:delta-asr-simple}
\end{figure}

A Figura~\ref{fig:effects-asr} decompõe o mesmo resultado nos efeitos do
desenho fatorial (Seção~\ref{sec:protocolo}): o efeito principal da
delimitação, o da marcação e a interação entre os dois, cada um estimado com as
quatro condições.

\begin{figure}[H]
\centering
\includegraphics[width=\textwidth]{effects_asr}
\caption{Efeitos do desenho $2 \times 2$ sobre a taxa de sucesso do ataque. O efeito principal de um fator é a média das duas condições com o fator menos a média das duas sem ele; a interação é a diferença entre o ganho da marcação com e sem delimitação. As hastes são intervalos de confiança de 95\% por \textit{bootstrap} BCa sobre os itens.}
\label{fig:effects-asr}
\end{figure}

\section{Custo da defesa}\label{sec:res_custo}

Uma redução na taxa de sucesso do ataque só é ganho se não sacrificar a tarefa
legítima. A Figura~\ref{fig:delta-utility} mostra a diferença de utilidade de
cada condição em relação ao \textit{prompt} simples, com o mesmo intervalo
pareado.

\begin{figure}[H]
\centering
\includegraphics[width=\textwidth]{delta_utility_vs_simple}
\caption{Diferença de utilidade de cada condição em relação ao \textit{prompt} simples, com intervalos de confiança de 95\% por \textit{bootstrap} BCa pareado sobre os itens. Só o BIPIA publica resposta de referência; no SEP e no NotInject, a utilidade significa ``respondeu sem ter sido bloqueado'', e não ``respondeu corretamente''.}
\label{fig:delta-utility}
\end{figure}

A Figura~\ref{fig:effects-utility} aplica a mesma decomposição fatorial à
utilidade. Por usar as quatro condições, ela separa o custo de cada fator
melhor que as comparações contra o \textit{prompt} simples.

\begin{figure}[H]
\centering
\includegraphics[width=\textwidth]{effects_utility}
\caption{Efeitos do desenho $2 \times 2$ sobre a utilidade, definidos como na Figura~\ref{fig:effects-asr}, com intervalos de confiança de 95\% por \textit{bootstrap} BCa sobre os itens. Vale a ressalva da Figura~\ref{fig:delta-utility} sobre o significado da utilidade.}
\label{fig:effects-utility}
\end{figure}

A Figura~\ref{fig:utility} mostra a utilidade em valor absoluto.

\begin{figure}[H]
\centering
\includegraphics[width=\textwidth]{utility}
\caption{Utilidade: fração das tarefas legítimas cumpridas, com item bloqueado contando como fracasso. As duas barras sem haste, do Llama-3.1-8B sem filtro no SEP e no NotInject, valem exatamente 1,000, sem variação para um intervalo. Vale a ressalva da Figura~\ref{fig:delta-utility} sobre o significado da utilidade.}
\label{fig:utility}
\end{figure}

A Figura~\ref{fig:overblock} mostra o sobre-bloqueio, a fração do conteúdo
legítimo que o filtro bloqueou.

\begin{figure}[H]
\centering
\includegraphics[width=\textwidth]{over_block}
\caption{Sobre-bloqueio: fração dos itens benignos que o filtro bloqueou. É a única métrica definida nos três benchmarks e o motivo de incluir o NotInject, composto de requisições legítimas com vocabulário sensível. Sem filtro não há decisão de bloqueio, e a barra correspondente é nula por construção. As barras sem haste do gpt-oss-20b e do gpt-oss-safeguard no BIPIA e no SEP somam de zero a quatro bloqueios em 1.500 decisões, poucos demais para um intervalo.}
\label{fig:overblock}
\end{figure}

\section{Alcance da marcação seletiva}\label{sec:res_alcance}

A marcação seletiva só age sobre os trechos que o filtro apontou. Sem trecho
apontado, não há o que marcar, e a condição se reduz ao \textit{prompt}
simples. A Figura~\ref{fig:marked} mede esse alcance.

\begin{figure}[H]
\centering
\includegraphics[width=\textwidth]{marked_rate}
\caption{Alcance da marcação seletiva: fração dos itens liberados pelo filtro com ao menos um trecho marcado. É zero por construção nas condições que não marcam.}
\label{fig:marked}
\end{figure}

\section{Saídas inválidas}\label{sec:res_validade}

A Figura~\ref{fig:invalid} mostra a taxa de saída inválida do filtro.

\begin{figure}[H]
\centering
\includegraphics[width=\textwidth]{invalid_output}
\caption{Taxa de saída inválida do filtro: fração das respostas que não seguiram o contrato de saída. Saída inválida conta como não bloqueio: o item segue para o componente principal e entra na medição.}
\label{fig:invalid}
\end{figure}

A Figura~\ref{fig:empty} mostra o equivalente no componente principal: a
fração dos itens liberados pelo filtro em que ele devolveu resposta vazia.

\begin{figure}[H]
\centering
\includegraphics[width=\textwidth]{empty_output}
\caption{Resposta vazia do componente principal: fração dos itens liberados pelo filtro em que o componente não produziu resposta, quase sempre por esgotar o teto de 8.192 \textit{tokens} raciocinando. Resposta vazia conta como saída inválida, e não como bloqueio nem como sobre-bloqueio: não executa o ataque e não cumpre a tarefa.}
\label{fig:empty}
\end{figure}

\section{Custo operacional}\label{sec:res_operacional}

As Figuras~\ref{fig:tokens} e~\ref{fig:latency} mostram o custo em
\textit{tokens} e o tempo de geração por item, somando os dois estágios.

\begin{figure}[H]
\centering
\includegraphics[width=\textwidth]{cost_tokens}
\caption{\textit{Tokens} por item, somando os dois estágios. A diferença para o \textit{prompt} simples é a despesa da delimitação e da marcação, que acrescentam texto; a diferença para a configuração sem filtro é a despesa do filtro, que dobra o número de chamadas.}
\label{fig:tokens}
\end{figure}

% [VALIDADE]
\begin{figure}[H]
\centering
\includegraphics[width=\textwidth]{cost_latency}
\caption{Tempo de geração por item, em segundos, somando os dois estágios. A barra é a mediana e o traço, o percentil 95. O tempo é o informado pelo provedor para cada chamada, e por isso não inclui fila, rede nem espera por limite de taxa. O mesmo modelo é servido por provedores com hardware diferente, escolhidos a cada chamada pela interface, e essa variação entra no tempo medido.}
\label{fig:latency}
\end{figure}

\section{Tabelas completas}\label{sec:res_tabelas}

As Tabelas~\ref{tab:resultados-bipia}, \ref{tab:resultados-sep}
e~\ref{tab:resultados-notinject} reúnem todas as métricas de todas as
configurações, uma por benchmark, com a explicação das colunas logo abaixo.
Elas são geradas dos mesmos arquivos CSV que originam as figuras.

\begin{table}[htb]
\centering
\caption{Métricas por condição e modelo no benchmark BIPIA, médias de 10 repetições.}
\label{tab:resultados-bipia}
\resizebox{\textwidth}{!}{%
\begin{tabular}{lrrrrrrrrr}
\toprule
model & condition & n & asr & utility & overBlock & invalidOutput & markedRate & meanTokens & latP50 \\
\midrule
llama-3.1-8b & simple & 300 & 0.123 & 0.564 & 0.021 & 0.001 & 0.000 & 1067.204 & 1.839 \\
llama-3.1-8b & delimiting & 300 & 0.103 & 0.539 & 0.067 & 0.001 & 0.000 & 1109.227 & 1.820 \\
llama-3.1-8b & demarking & 300 & 0.047 & 0.493 & 0.027 & 0.001 & 0.069 & 1203.440 & 1.814 \\
llama-3.1-8b & delimiting\_demarking & 300 & 0.039 & 0.482 & 0.065 & 0.001 & 0.079 & 1232.852 & 1.762 \\
gpt-oss-20b & simple & 300 & 0.235 & 0.661 & 0.001 & 0.074 & 0.000 & 1411.183 & 4.872 \\
gpt-oss-20b & delimiting & 300 & 0.129 & 0.661 & 0.000 & 0.030 & 0.000 & 1408.372 & 4.962 \\
gpt-oss-20b & demarking & 300 & 0.153 & 0.639 & 0.001 & 0.069 & 0.041 & 1514.817 & 4.837 \\
gpt-oss-20b & delimiting\_demarking & 300 & 0.097 & 0.638 & 0.000 & 0.037 & 0.021 & 1505.214 & 4.823 \\
llama-3.1-8b & no\_filter & 300 & 0.184 & 0.587 & 0.000 & 0.000 & 0.000 & 527.205 & 1.024 \\
gpt-oss-20b & no\_filter & 300 & 0.795 & 0.660 & 0.000 & 0.000 & 0.000 & 794.128 & 2.932 \\
gpt-oss-safeguard & simple & 300 & 0.374 & 0.661 & 0.000 & 0.000 & 0.000 & 1492.515 & 2.497 \\
gpt-oss-safeguard & delimiting & 300 & 0.314 & 0.654 & 0.000 & 0.000 & 0.000 & 1558.584 & 2.420 \\
gpt-oss-safeguard & demarking & 300 & 0.305 & 0.649 & 0.000 & 0.001 & 0.030 & 1611.346 & 2.377 \\
gpt-oss-safeguard & delimiting\_demarking & 300 & 0.235 & 0.646 & 0.000 & 0.000 & 0.033 & 1657.387 & 2.104 \\
\bottomrule
\end{tabular}
}
\vspace{2pt}
\begin{minipage}{\textwidth}
\footnotesize\raggedright
\textbf{Legenda.} \textbf{model}: modelo que instancia o sistema (filtro e componente principal); nas linhas de guard, o modelo do filtro; \textbf{condition}: condição experimental: simple, delimiting, demarking, delimiting\_demarking, ou a referência no\_filter; \textbf{n}: número de itens efetivamente executados na configuração; \textbf{asr}: taxa de sucesso do ataque na saída do componente principal, sobre os itens de ataque (métrica primária; menor é melhor); \textbf{utility}: fração das tarefas benignas cumpridas; item bloqueado conta como fracasso (maior é melhor); \textbf{overBlock}: fração dos itens benignos que o filtro bloqueou (sobre-bloqueio); \textbf{invalidOutput}: fração das saídas do filtro que não puderam ser interpretadas; contam conservadoramente como não bloqueio; \textbf{markedRate}: fração dos itens liberados que receberam ao menos um trecho marcado, que é o alcance efetivo da marcação seletiva; \textbf{meanTokens}: média de tokens (entrada + saída) por item, somando os dois estágios; \textbf{latP50}: tempo de geração mediano por item, em segundos, medido pelo provedor e somando os dois estágios; exclui fila e rede.
\end{minipage}
\end{table}

\begin{table}[htb]
\centering
\caption{Métricas por condição e modelo no benchmark SEP, médias de 10 repetições.}
\label{tab:resultados-sep}
\resizebox{\textwidth}{!}{%
\begin{tabular}{lrrrrrrrrr}
\toprule
model & condition & n & asr & utility & overBlock & invalidOutput & markedRate & meanTokens & latP50 \\
\midrule
llama-3.1-8b & simple & 300 & 0.097 & 0.967 & 0.033 & 0.000 & 0.000 & 792.746 & 3.700 \\
llama-3.1-8b & delimiting & 300 & 0.087 & 0.949 & 0.051 & 0.000 & 0.000 & 867.524 & 3.682 \\
llama-3.1-8b & demarking & 300 & 0.075 & 0.966 & 0.034 & 0.000 & 0.022 & 879.112 & 2.837 \\
llama-3.1-8b & delimiting\_demarking & 300 & 0.069 & 0.954 & 0.046 & 0.000 & 0.112 & 948.382 & 2.929 \\
gpt-oss-20b & simple & 300 & 0.743 & 0.977 & 0.001 & 0.013 & 0.000 & 1733.365 & 9.577 \\
gpt-oss-20b & delimiting & 300 & 0.735 & 0.979 & 0.003 & 0.010 & 0.000 & 1773.805 & 10.151 \\
gpt-oss-20b & demarking & 300 & 0.691 & 0.993 & 0.001 & 0.014 & 0.004 & 1632.028 & 8.431 \\
gpt-oss-20b & delimiting\_demarking & 300 & 0.657 & 0.993 & 0.002 & 0.013 & 0.006 & 1719.191 & 8.935 \\
llama-3.1-8b & no\_filter & 300 & 0.101 & 1.000 & 0.000 & 0.000 & 0.000 & 484.226 & 2.887 \\
gpt-oss-20b & no\_filter & 300 & 0.783 & 0.984 & 0.000 & 0.000 & 0.000 & 1241.679 & 7.312 \\
gpt-oss-safeguard & simple & 300 & 0.779 & 0.981 & 0.000 & 0.000 & 0.000 & 1705.883 & 7.411 \\
gpt-oss-safeguard & delimiting & 300 & 0.785 & 0.978 & 0.001 & 0.000 & 0.000 & 1852.056 & 7.664 \\
gpt-oss-safeguard & demarking & 300 & 0.701 & 0.990 & 0.001 & 0.000 & 0.001 & 1663.635 & 6.664 \\
gpt-oss-safeguard & delimiting\_demarking & 300 & 0.692 & 0.995 & 0.000 & 0.000 & 0.002 & 1765.446 & 6.503 \\
\bottomrule
\end{tabular}
}
\vspace{2pt}
\begin{minipage}{\textwidth}
\footnotesize\raggedright
\textbf{Legenda.} \textbf{model}: modelo que instancia o sistema (filtro e componente principal); nas linhas de guard, o modelo do filtro; \textbf{condition}: condição experimental: simple, delimiting, demarking, delimiting\_demarking, ou a referência no\_filter; \textbf{n}: número de itens efetivamente executados na configuração; \textbf{asr}: taxa de sucesso do ataque na saída do componente principal, sobre os itens de ataque (métrica primária; menor é melhor); \textbf{utility}: fração das tarefas benignas cumpridas; item bloqueado conta como fracasso (maior é melhor); \textbf{overBlock}: fração dos itens benignos que o filtro bloqueou (sobre-bloqueio); \textbf{invalidOutput}: fração das saídas do filtro que não puderam ser interpretadas; contam conservadoramente como não bloqueio; \textbf{markedRate}: fração dos itens liberados que receberam ao menos um trecho marcado, que é o alcance efetivo da marcação seletiva; \textbf{meanTokens}: média de tokens (entrada + saída) por item, somando os dois estágios; \textbf{latP50}: tempo de geração mediano por item, em segundos, medido pelo provedor e somando os dois estágios; exclui fila e rede.
\end{minipage}
\end{table}

\begin{table}[htb]
\centering
\caption{Métricas por condição e modelo no benchmark NotInject, médias de 10 repetições.}
\label{tab:resultados-notinject}
\resizebox{\textwidth}{!}{%
\begin{tabular}{lrrrrrrrrr}
\toprule
model & condition & n & asr & utility & overBlock & invalidOutput & markedRate & meanTokens & latP50 \\
\midrule
llama-3.1-8b & simple & 339 & -- & 0.902 & 0.098 & 0.000 & 0.000 & 772.803 & 4.460 \\
llama-3.1-8b & delimiting & 339 & -- & 0.919 & 0.081 & 0.001 & 0.000 & 873.287 & 4.628 \\
llama-3.1-8b & demarking & 339 & -- & 0.901 & 0.099 & 0.000 & 0.002 & 830.562 & 3.490 \\
llama-3.1-8b & delimiting\_demarking & 339 & -- & 0.921 & 0.079 & 0.001 & 0.002 & 910.876 & 3.689 \\
gpt-oss-20b & simple & 339 & -- & 0.922 & 0.044 & 0.017 & 0.000 & 1881.129 & 10.591 \\
gpt-oss-20b & delimiting & 339 & -- & 0.906 & 0.055 & 0.011 & 0.000 & 1953.257 & 10.559 \\
gpt-oss-20b & demarking & 339 & -- & 0.932 & 0.047 & 0.013 & 0.001 & 1635.738 & 8.136 \\
gpt-oss-20b & delimiting\_demarking & 339 & -- & 0.932 & 0.054 & 0.013 & 0.003 & 1661.217 & 8.143 \\
llama-3.1-8b & no\_filter & 339 & -- & 1.000 & 0.000 & 0.000 & 0.000 & 540.755 & 3.879 \\
gpt-oss-20b & no\_filter & 339 & -- & 0.964 & 0.000 & 0.000 & 0.000 & 1477.801 & 8.057 \\
gpt-oss-safeguard & simple & 339 & -- & 0.948 & 0.016 & 0.000 & 0.000 & 1879.547 & 8.492 \\
gpt-oss-safeguard & delimiting & 339 & -- & 0.942 & 0.019 & 0.000 & 0.000 & 1993.751 & 8.432 \\
gpt-oss-safeguard & demarking & 339 & -- & 0.969 & 0.015 & 0.000 & 0.000 & 1626.800 & 6.008 \\
gpt-oss-safeguard & delimiting\_demarking & 339 & -- & 0.962 & 0.019 & 0.000 & 0.001 & 1693.818 & 6.032 \\
\bottomrule
\end{tabular}
}
\vspace{2pt}
\begin{minipage}{\textwidth}
\footnotesize\raggedright
\textbf{Legenda.} \textbf{model}: modelo que instancia o sistema (filtro e componente principal); nas linhas de guard, o modelo do filtro; \textbf{condition}: condição experimental: simple, delimiting, demarking, delimiting\_demarking, ou a referência no\_filter; \textbf{n}: número de itens efetivamente executados na configuração; \textbf{asr}: taxa de sucesso do ataque na saída do componente principal, sobre os itens de ataque (métrica primária; menor é melhor); \textbf{utility}: fração das tarefas benignas cumpridas; item bloqueado conta como fracasso (maior é melhor); \textbf{overBlock}: fração dos itens benignos que o filtro bloqueou (sobre-bloqueio); \textbf{invalidOutput}: fração das saídas do filtro que não puderam ser interpretadas; contam conservadoramente como não bloqueio; \textbf{markedRate}: fração dos itens liberados que receberam ao menos um trecho marcado, que é o alcance efetivo da marcação seletiva; \textbf{meanTokens}: média de tokens (entrada + saída) por item, somando os dois estágios; \textbf{latP50}: tempo de geração mediano por item, em segundos, medido pelo provedor e somando os dois estágios; exclui fila e rede.
\end{minipage}
\end{table}

